
{\centering \section*{附件A: 支撑材料文件列表}}

\begin{table}[htbp]
  \centering
  \begin{tabular}{ll}
    \toprule
    文件名 & 简介 \\
    \midrule
    % Python代码文件部分
    数据预处理.py & 数据预处理相关代码实现 \\
    问题1\_相关分析.py & 问题1的相关分析代码 \\
    问题1\_LMM.py & 问题1的LMM模型实现代码 \\
    问题1\_GAMM.py & 问题1的GAMM模型实现代码 \\
    问题2.py & 问题2的求解代码 \\
    问题3.py & 问题3的求解代码 \\
    问题4.py & 问题4的求解代码 \\
    \midrule
    % 数据与结果文件部分
    AI 工具使用详情说明.pdf & 本次数模比赛AI工具使用报告 \\
    GAMM\_coef\_table\_weighted.xlsx & 问题1 GAMM模型各指标信息 \\
    GAMM\_model\_summary\_weighted.txt & 问题1 GAMM建模过程信息报告 \\
    聚类自助法\_斯皮尔曼相关分析结果\_带权.xlsx & 问题1 斯皮尔曼相关分析表格 \\
    聚类自助法\_所有配对结果\_带权.xlsx & 问题1 所有配对结果表 \\
    聚类自助法\_显著配对结果\_带权.xlsx & 问题1 显著配对结果表 \\
    bic\_table.txt & 问题2 BMI最优分组决策数据 \\
    bmi\_groups\_results.csv & 问题2 最优分组的各组信息 \\
    run\_summary.txt & 问题2 建模过程信息报告 \\
    Ti\_table.csv & 问题2 个体首次达标时间信息表 \\
    groups\_policy\_improved.csv & 问题3 BMI最优分组的各组信息 \\
    summary\_improved.txt & 问题3 建模过程信息报告 \\
    best\_model.pkl & 问题4 训练的最佳性能机器学习模型 \\
    best\_model\_info.json & 问题4 记录最佳模型的参数 \\
    analysis\_report.txt & 问题4 机器学习模型训练报告 \\
    model\_metrics.xlsx & 问题4 各模型多性能指标汇总 \\
    resample\_stats.xlsx & 问题4 训练各模型的数据集特征汇总 \\
    processed\_data.xlsx & 问题4 处理过程中对附件数据的加工 \\
    \bottomrule
  \end{tabular}
  \label{tab:file清单}
\end{table}

{\centering \section*{附录B: 核心python代码}}

\noindent \textbf{0.1} 这是数据预处理所用核心代码，进行数据读取并标准化，计算质量阈值，软门控判定并对样本质量评分，为进一步分析做好准备。
% 第一个附录 appendices是附录环境
\begin{appendices}
\begin{lstlisting}

###数据预处理_核心代码


#软门控判定
def qc_soft_flags(row, thr):
    flags = []
    O  = pd.to_numeric(row.get(COL_O),  errors="coerce")
    M  = pd.to_numeric(row.get(COL_M),  errors="coerce")
    N  = pd.to_numeric(row.get(COL_N),  errors="coerce")
    AA = pd.to_numeric(row.get(COL_AA), errors="coerce")
    P  = pd.to_numeric(row.get(COL_P),  errors="coerce")
    GC13 = pd.to_numeric(row.get(COL_GC13), errors='coerce')
    GC18 = pd.to_numeric(row.get(COL_GC18), errors='coerce')
    GC21 = pd.to_numeric(row.get(COL_GC21), errors='coerce')

    # 比率/计数
    if (not np.isfinite(M)) or (M < thr['M5']):
        flags.append("low_align(<M5)")
    if (not np.isfinite(N)) or (N > thr['N95']):
        flags.append("high_dup(>N95)")
    if (not np.isfinite(AA)) or (AA > thr['AA95']):
        flags.append("high_filtered(>AA95)")
    if (not np.isfinite(O)) or (O < thr['O5']):
        flags.append("low_unique(<O5)")

    # 整体 GC
    if (not np.isfinite(P)) or (P < thr['GC_LOW']):
        flags.append("gc_overall_low(<P5)")
    if np.isfinite(P) and (P > thr['GC_HIGH']):
        flags.append("gc_overall_high(>0.60)")

    # ΔGC
    if np.isfinite(GC13) and np.isfinite(GC18) and np.isfinite(GC21) and np.isfinite(P):
        for k, gck in [('13',GC13),('18',GC18),('21',GC21)]:
            z = ((gck - P) - thr['mu_d'][k]) / (thr['sd_d'][k] + 1e-9)
            if abs(z) > 3:
                flags.append(f"deltaGC{k}_>|3sd|")
                break
    else:
        flags.append("deltaGC_missing")

    return flags

# 对被软门控样本计算 Q 
def compute_Q_for_flagged(row, thr, return_parts=False):
    """返回 Q（及可选的各组成分）"""
    def clip01(x): return np.minimum(1.0, np.maximum(0.0, x))
    O  = pd.to_numeric(row.get(COL_O),  errors="coerce")
    M  = pd.to_numeric(row.get(COL_M),  errors="coerce")
    N  = pd.to_numeric(row.get(COL_N),  errors="coerce")
    AA = pd.to_numeric(row.get(COL_AA), errors="coerce")
    P  = pd.to_numeric(row.get(COL_P),  errors="coerce")
    GC13 = pd.to_numeric(row.get(COL_GC13), errors='coerce')
    GC18 = pd.to_numeric(row.get(COL_GC18), errors='coerce')
    GC21 = pd.to_numeric(row.get(COL_GC21), errors='coerce')

    # 分位区间（缺失兜底）
    O5 = thr['O5']; O95 = O5 * 1.5 if np.isfinite(O5) else (O if np.isfinite(O) else 1.0)
    M5 = thr['M5']; M95 = min(0.95, M5 + 0.15) if np.isfinite(M5) else 0.90
    N5, N95 = 0.0, (thr['N95'] if np.isfinite(thr['N95']) else 0.04)
    AA5, AA95 = 0.0, (thr['AA95'] if np.isfinite(thr['AA95']) else 0.04)
    eps = 1e-9

    f_reads = clip01((np.log10(max(O,1.0)) - np.log10(max(O5,1.0))) /
                     (np.log10(max(O95, O5+1)) - np.log10(max(O5,1.0)) + eps)) if np.isfinite(O) and np.isfinite(O5) else 0.5
    f_align = clip01((M - M5) / (M95 - M5 + eps)) if np.isfinite(M) and np.isfinite(M5) else 0.5
    f_dup   = 1.0 - clip01((N - N5) / (N95 - N5 + eps)) if np.isfinite(N) and np.isfinite(N95) else 0.5
    f_filt  = 1.0 - clip01((AA - AA5)/ (AA95- AA5+ eps)) if np.isfinite(AA) and np.isfinite(AA95) else 0.5

    # 整体 GC：温和惩罚（0.6~1.0）
    f_gc_overall = 1.0
    if np.isfinite(P):
        if P < thr['GC_LOW']:
            P1 = thr['P1'] if np.isfinite(thr['P1']) else max(thr['GC_LOW'] - 0.02, 0.0)
            span = max(thr['GC_LOW'] - P1, 1e-3)
            frac = (P - P1) / span  # [0,1]
            f_gc_overall = 0.6 + 0.4 * clip01(frac)
        elif P > thr['GC_HIGH']:
            P99 = thr['P99'] if np.isfinite(thr['P99']) else (thr['GC_HIGH'] + 0.02)
            span = max(P99 - thr['GC_HIGH'], 1e-3)
            frac = (P99 - P) / span
            f_gc_overall = 0.6 + 0.4 * clip01(frac)
        else:
            f_gc_overall = 1.0

    # ΔGC：取 13/18/21 的最小值
    f_delta = 1.0
    if np.isfinite(GC13) and np.isfinite(GC18) and np.isfinite(GC21) and np.isfinite(P):
        f_list = []
        for k, gck in [('13',GC13),('18',GC18),('21',GC21)]:
            z = ((gck - P) - thr['mu_d'][k]) / (thr['sd_d'][k] + eps)
            f_list.append(1.0 - min(1.0, abs(z)/3.0))
        f_delta = min(f_list)

    # 几何平均（读段/比对加倍权重）
    prod = (f_reads**2) * (f_align**2) * f_dup * f_filt * f_gc_overall * f_delta
    Q = float(max(prod ** (1/8), Q_FLOOR))

    if return_parts:
        return Q, dict(f_reads=f_reads, f_align=f_align, f_dup=f_dup, f_filt=f_filt,
                       f_gc_overall=f_gc_overall, f_delta=f_delta)
    return Q

# 选“质量最好”的记录（同孕妇同一天仅保留一条）
def pick_best_per_day(df):
    # ΔGC 最大绝对偏差（tie-breaker）
    d13 = pd.to_numeric(df[COL_GC13], errors='coerce') - pd.to_numeric(df[COL_P], errors='coerce')
    d18 = pd.to_numeric(df[COL_GC18], errors='coerce') - pd.to_numeric(df[COL_P], errors='coerce')
    d21 = pd.to_numeric(df[COL_GC21], errors='coerce') - pd.to_numeric(df[COL_P], errors='coerce')
    df["_deltaGC_maxabs"] = pd.concat([d13.abs(), d18.abs(), d21.abs()], axis=1).max(axis=1)

\end{lstlisting}
\end{appendices}

\noindent \textbf{0.2} 这是问题一相关性分析所用核心代码，基于聚类自助法计算了带权相关性考虑了数据的聚类，并评估相关性的显著性，提供p值。
\begin{appendices}
\begin{lstlisting}

import pandas as pd
import numpy as np
from scipy import stats
import matplotlib.pyplot as plt
import seaborn as sns
import warnings, time
from tqdm import tqdm
warnings.filterwarnings('ignore')
plt.rcParams['font.sans-serif'] = ['SimHei']
plt.rcParams['axes.unicode_minus'] = False


# 带权Spearman：秩→加权Pearson
def weighted_spearmanr(x, y, w):
    x = np.asarray(x, dtype=float)
    y = np.asarray(y, dtype=float)
    w = np.asarray(w, dtype=float)

    mask = np.isfinite(x) & np.isfinite(y) & np.isfinite(w) & (w > 0)
    if mask.sum() <= 2:
        return np.nan

    xr = stats.rankdata(x[mask], method='average')
    yr = stats.rankdata(y[mask], method='average')
    ww = w[mask]

    ww_sum = ww.sum()
    if ww_sum <= 0:
        return np.nan

    def wmean(arr, ww): return (ww * arr).sum() / ww_sum
    xbar = wmean(xr, ww)
    ybar = wmean(yr, ww)
    xc = xr - xbar
    yc = yr - ybar

    cov_w = (ww * xc * yc).sum() / ww_sum
    varx = (ww * xc * xc).sum() / ww_sum
    vary = (ww * yc * yc).sum() / ww_sum
    if varx <= 0 or vary <= 0:
        return np.nan
    return cov_w / np.sqrt(varx * vary)

#聚类自助法+带权Spearman
def cluster_bootstrap_weighted_spearman(
    data, col1, col2, cluster_col, weight_col,
    n_bootstrap=1000, confidence_level=0.95, random_state=None
):
    valid = ~(data[col1].isna() | data[col2].isna())
    if weight_col in data.columns:
        w = pd.to_numeric(data[weight_col], errors='coerce').fillna(0.0)
    else:
        w = pd.Series(1.0, index=data.index)

    original_corr = weighted_spearmanr(data.loc[valid, col1], data.loc[valid, col2], w.loc[valid])

    uniq = data[cluster_col].astype(str).unique()
    n_clusters = len(uniq)
    cluster_dict = {}
    for c in uniq:
        sub = data[data[cluster_col].astype(str) == str(c)][[col1, col2]].copy()
        sub['__w__'] = w.loc[sub.index].values
        cluster_dict[c] = sub.values  # shape: (n_i, 3)

    rng = np.random.default_rng(random_state)
    boot_corrs = np.zeros(n_bootstrap, dtype=float)
    valid_count = 0

    for i in range(n_bootstrap):
        # 有放回抽簇
        picked = rng.choice(uniq, size=n_clusters, replace=True)
        parts = [cluster_dict[c] for c in picked if c in cluster_dict]
        if not parts:
            continue
        samp = np.vstack(parts)  # [:,0]=col1, [:,1]=col2, [:,2]=w
        x, y, ww = samp[:, 0], samp[:, 1], samp[:, 2]
        corr = weighted_spearmanr(x, y, ww)
        if np.isfinite(corr):
            boot_corrs[valid_count] = corr
            valid_count += 1

    if valid_count == 0:
        return dict(correlation=original_corr, p_value=np.nan, ci_lower=np.nan,
                    ci_upper=np.nan, se=np.nan, n_bootstrap=0)

    boot_corrs = boot_corrs[:valid_count]
    se = float(np.std(boot_corrs))
    alpha = 1 - confidence_level
    ci_lower = float(np.percentile(boot_corrs, alpha/2 * 100))
    ci_upper = float(np.percentile(boot_corrs, (1 - alpha/2) * 100))

    # 简单双尾p值：看自助分布相对于0的尾部概率
    if original_corr >= 0:
        p_value = 2 * np.mean(boot_corrs <= 0)
    else:
        p_value = 2 * np.mean(boot_corrs >= 0)
    p_value = float(min(1.0, max(0.0, p_value)))

    return dict(
        correlation=float(original_corr),
        p_value=p_value,
        ci_lower=ci_lower,
        ci_upper=ci_upper,
        se=se,
        n_bootstrap=int(valid_count)
    )


def main():
    df, used_path, used_sheet ,last_err= None, None, None,None
    df = pd.read_excel("附件-男胎检测数据-处理.xlsx")

    #权重列
    weight_col = 'w_Q'
    if weight_col not in df.columns:
        print("警告：未找到 w_Q 列，默认权重=1")
        df[weight_col] = 1.0
    df[weight_col] = pd.to_numeric(df[weight_col], errors='coerce').fillna(1.0).clip(lower=0)

    columns_to_analyze = [
        '年龄','检测孕周','孕妇BMI','原始读段数','在参考基因组上比对的比例','重复读段的比例',
        '唯一比对的读段数','GC含量','13号染色体的Z值','18号染色体的Z值','21号染色体的Z值',
        'X染色体的Z值','Y染色体的Z值','Y染色体浓度','X染色体浓度',
        '13号染色体的GC含量','18号染色体的GC含量','21号染色体的GC含量','被过滤掉读段数的比例'
    ]
    exists = [c for c in columns_to_analyze if c in df.columns]
    if len(exists) < len(columns_to_analyze):
        missing = sorted(set(columns_to_analyze) - set(exists))
        print(f"提示：以下列不存在，将自动忽略：{missing}")
    columns_to_analyze = exists
    
    #聚类列
    cluster_col = '孕妇代码'
    if cluster_col not in df.columns:
        for cand in ['孕妇ID','孕妇编号','ID','编号','代码']:
            if cand in df.columns:
                cluster_col = cand
                break
        else:
            print("未找到聚类标识列，将使用索引创建虚拟簇")
            df[cluster_col] = df.index

    #只取需要的列 + 转数值
    data = df[columns_to_analyze + [cluster_col, weight_col]].copy()
    for c in columns_to_analyze:
        data[c] = pd.to_numeric(data[c], errors='coerce')

    uniq = data[cluster_col].nunique()
    total = len(data)
    print(f"独立孕妇数: {uniq}；总观察数: {total}；平均每孕妇检测次数: {total/uniq:.2f}")
    print("缺失值统计：")
    print(data[columns_to_analyze].isnull().sum())

    n_bootstrap = 1000
    print(f"\n开始带权聚类自助法（B={n_bootstrap}）…\n")

    n_vars = len(columns_to_analyze)
    corrM = np.zeros((n_vars, n_vars)); pM = np.zeros_like(corrM)
    loM = np.zeros_like(corrM); upM = np.zeros_like(corrM)
    seM = np.zeros_like(corrM)

    start = time.time()
    tot_pairs = n_vars * (n_vars + 1) // 2
    with tqdm(total=tot_pairs, desc="计算进度") as pbar:
        for i, c1 in enumerate(columns_to_analyze):
            for j in range(i, n_vars):
                c2 = columns_to_analyze[j]
                if i == j:
                    corrM[i,j]=1.0; pM[i,j]=0.0; loM[i,j]=1.0
                    upM[i,j]=1.0; seM[i,j]=0.0
                else:
                    res = cluster_bootstrap_weighted_spearman(
                        data, c1, c2, cluster_col, weight_col,
                        n_bootstrap=n_bootstrap, confidence_level=0.95
                    )
                    corrM[i,j]=corrM[j,i]=res['correlation']
                    pM[i,j]=pM[j,i]=res['p_value']
                    loM[i,j]=loM[j,i]=res['ci_lower']
                    upM[i,j]=upM[j,i]=res['ci_upper']
                    seM[i,j]=seM[j,i]=res['se']
                pbar.update(1)
    dur = time.time() - start
    print(f"\n完成！用时 {dur:.1f}s")

    #输出与可视化
    spearman_df = pd.DataFrame(corrM, index=columns_to_analyze, 
                               columns=columns_to_analyze)
    p_df = pd.DataFrame(pM,    index=columns_to_analyze, 
                               columns=columns_to_analyze)
    lo_df = pd.DataFrame(loM,   index=columns_to_analyze, 
                               columns=columns_to_analyze)
    up_df = pd.DataFrame(upM,   index=columns_to_analyze, 
                               columns=columns_to_analyze)
    se_df = pd.DataFrame(seM,   index=columns_to_analyze, 
                               columns=columns_to_analyze)

    sig_df = pd.DataFrame(index=columns_to_analyze, columns=columns_to_analyze)
    for i in range(n_vars):
        for j in range(n_vars):
            if i==j: sig_df.iloc[i,j]='-'; continue
            pv=pM[i,j]
            if np.isnan(pv): sig='NA'
            elif pv<1e-3: sig='***'
            elif pv<1e-2: sig='**'
            elif pv<5e-2: sig='*'
            else: sig='ns'
            sig_df.iloc[i,j]=sig

    #保存
    with pd.ExcelWriter('聚类自助法_斯皮尔曼相关分析结果_带权.xlsx', engine='openpyxl') as w:
        info = pd.DataFrame({
            '项目':['分析方法','是否使用权重',
                  '权重列','独立孕妇数','总观察数',
                  'Bootstrap次数','置信水平','计算用时(秒)',
                  '数据文件','工作表'],
            '值'  :['Cluster Bootstrap + Weighted Spearman',
                   '是','w_Q',uniq,total,n_bootstrap,'95%',
                   f'{dur:.1f}',used_path,used_sheet]
        })
        info.to_excel(w, sheet_name='分析信息', index=False)
        spearman_df.to_excel(w, sheet_name='相关系数矩阵')
        p_df.to_excel(w,sheet_name='P值矩阵_聚类校正(带权)')
        lo_df.to_excel(w,sheet_name='置信区间下限')
        up_df.to_excel(w,sheet_name='置信区间上限')
        se_df.to_excel(w,sheet_name='标准误')
        sig_df.to_excel(w,sheet_name='显著性标记')
        data[columns_to_analyze].describe().to_excel(w, sheet_name='描述统计(未加权)')

    #配对列表
    rows=[]
    for i in range(n_vars):
        for j in range(i+1, n_vars):
            rows.append({
                '指标1':columns_to_analyze[i],
                '指标2':columns_to_analyze[j],
                '相关系数':corrM[i,j],
                'P值(聚类校正,带权)':pM[i,j],
                '95%下限':loM[i,j],
                '95%上限':upM[i,j],
                '标准误':seM[i,j],
                '显著性':sig_df.iloc[i,j],
                '置信区间包含0': '是' if (np.isfinite(loM[i,j]) and loM[i,j] <= 0 <= upM[i,j]) 
                else '否'
            })
    pairs = pd.DataFrame(rows).sort_values(['P值(聚类校正,带权)','相关系数'], 
                                           ascending=[True, False])
    pairs.to_excel('聚类自助法_所有配对结果_带权.xlsx', index=False)
    sig_pairs = pairs[pairs['P值(聚类校正,带权)']<0.05]
    if not sig_pairs.empty:
        sig_pairs.to_excel('聚类自助法_显著配对结果_带权.xlsx', index=False)

    # 可视化
    plt.figure(figsize=(18,16))
    sns.heatmap(spearman_df, annot=True, fmt='.2f', 
                cmap='RdBu_r', center=0, vmin=-1, 
                vmax=1,square=True, linewidths=0.5, 
                cbar_kws={"shrink":0.8,"label":"带权Spearman相关"})
    plt.title('带权Spearman相关系数热力图（聚类自助法）', 
              fontsize=18, pad=20)
    plt.xticks(rotation=45, ha='right'); plt.yticks(rotation=0)
    plt.tight_layout()
    plt.savefig('聚类自助法_相关系数热力图_带权.png', dpi=300)

    plt.figure(figsize=(18,16))
    sns.heatmap(p_df, annot=True, fmt='.3f', cmap='YlOrRd_r',
                vmin=0, vmax=1,square=True, linewidths=0.5, 
                cbar_kws={"shrink":0.8,"label":"P值（聚类校正,带权）"})
    plt.title('P值热力图（聚类自助法校正, 带权）', fontsize=18, pad=20)
    plt.xticks(rotation=45, ha='right'); plt.yticks(rotation=0)
    plt.tight_layout(); plt.savefig('聚类自助法_P值热力图_带权.png', dpi=300)

    # 显著性热力图
    sig_numeric = np.zeros((n_vars, n_vars))
    mapping = {'***':3,'**':2,'*':1,'ns':0,'-':-1}
    for i in range(n_vars):
        for j in range(n_vars):
            sig_numeric[i,j] = mapping.get(sig_df.iloc[i,j], np.nan)
    plt.figure(figsize=(18,16))
    cmap = sns.color_palette(['#d4d4d4','#808080','#ffcccc','#ff6666','#ff0000'], 
                             as_cmap=True)
    sns.heatmap(sig_numeric, annot=sig_df.values, fmt='s', cmap=cmap, vmin=-1, vmax=3,
                square=True, linewidths=0.5, cbar_kws={"shrink":0.8,"label":"显著性水平"})
    cbar = plt.gcf().axes[-1]
    cbar.set_yticklabels(['对角线','ns','*','**','***'])
    plt.title('显著性热力图（聚类自助法校正, 带权）', fontsize=18, pad=20)
    plt.xticks(rotation=45, ha='right'); plt.yticks(rotation=0)
    plt.tight_layout(); plt.savefig('聚类自助法_显著性热力图_带权.png', dpi=300)

if __name__ == '__main__':
    main()



\end{lstlisting}
\end{appendices}

\noindent \textbf{0.3} 这是广义加性混合模型GAMM求解核心程序，主要目标是分析孕周、BMI等因素与Y染色体浓度之间的关系。它对数据进行预处理，处理缺失值，应用加权分析，并通过使用立方体回归样条和混合线性模型探索最佳的GAMM配置。
\begin{appendices}
\begin{lstlisting}

import re, warnings
import numpy as np
import pandas as pd
import statsmodels.api as sm
from scipy import stats
from patsy import dmatrix, build_design_matrices
import matplotlib, matplotlib.pyplot as plt
from pathlib import Path

matplotlib.rcParams["font.sans-serif"] = ["Microsoft YaHei", "SimHei", "Arial Unicode MS", "DejaVu Sans"]
matplotlib.rcParams["axes.unicode_minus"] = False

CANDIDATES = [
    ("附件-男胎检测数据-处理.xlsx", "附件-男胎检测数据-处理"),
    ("附件.xlsx", "男胎检测数据"),
]
COL_ID, COL_Y, COL_GA, COL_BMI, COL_AGE, COL_W = "孕妇代码", "Y染色体浓度", "检测孕周", "孕妇BMI", "年龄", "w_Q"

out_dir = Path("./model_outputs_gamm_w"); out_dir.mkdir(parents=True, exist_ok=True)

def parse_gestation_to_weeks(x):
    if pd.isna(x): return np.nan
    s = str(x).strip()
    try: return float(s)
    except: pass
    m = re.search(r'(\d+)\s*(?:周|w|W)?\s*\+?\s*(\d+)?\s*(?:天|d|D)?', s)
    if m:
        w = float(m.group(1)); d = float(m.group(2)) if m.group(2) else 0.0
        return w + d/7.0
    m2 = re.search(r'(\d+)', s)
    return float(m2.group(1)) if m2 else np.nan

def smithson_verkuilen(y):
    y = y.clip(0,1); n = y.notna().sum()
    return (y*(n-1)+0.5)/n

def inv_logit(z): return 1/(1+np.exp(-z))
def AIC(llf, k): return 2*k - 2*llf
def BIC(llf, k, n): return np.log(n)*k - 2*llf

def wmean(x, w):
    x = pd.to_numeric(x, errors="coerce"); w = pd.to_numeric(w, errors="coerce").clip(lower=0)
    m = np.nansum(w * x) / np.nansum(w) if np.nansum(w) > 0 else np.nan
    return m

def fit_mixedlm_weighted(endog, exog, groups, weights, exog_re=None):
    w = np.asarray(weights, dtype=float)
    try:
        md = sm.MixedLM(endog, exog, groups=groups, exog_re=exog_re, weights=w)
        return md.fit(reml=False, method="lbfgs")
    except (TypeError, ValueError):
        print("当前 statsmodels MixedLM 不支持 weights")
        sw = np.sqrt(np.clip(w, 1e-8, None))
        endog_t = endog * sw
        exog_t  = exog * sw[:, None]
        exre_t  = exog_re * sw[:, None] if exog_re is not None else None
        md = sm.MixedLM(endog_t, exog_t, groups=groups, exog_re=exre_t)
        return md.fit(reml=False, method="lbfgs")

def lrt(res_small, res_large, df_diff=None):
    LR = 2*(res_large.llf - res_small.llf)
    if df_diff is None:
        df_diff = len(res_large.params) - len(res_small.params)
    p = stats.chi2.sf(LR, df_diff)
    return LR, df_diff, p

#读取数据
last_err = None
for path, sheet in CANDIDATES:
    try:
        df = pd.read_excel(path, sheet_name=sheet)
        used_path, used_sheet = path, sheet
        break
    except Exception as e:
        last_err = e

for c in [COL_ID, COL_Y, COL_GA, COL_BMI]:
    if c not in df.columns:
        raise SystemExit(f"找不到列：{c}（数据源：{used_path}/{used_sheet}）")
if COL_W not in df.columns:
    print("[提示] 未找到 w_Q 列，默认权重=1")
    df[COL_W] = 1.0

use_cols = [COL_ID, COL_Y, COL_GA, COL_BMI, COL_W] + ([COL_AGE] if COL_AGE in df.columns else [])
df = df[use_cols].copy()
df.rename(columns={COL_ID:"id", COL_Y:"Y", COL_GA:"GA_raw", COL_BMI:"BMI_raw", COL_W:"w"}, inplace=True)
if COL_AGE in df.columns: df.rename(columns={COL_AGE:"Age_raw"}, inplace=True)

# 变量转换
df["GA"]  = df["GA_raw"].apply(parse_gestation_to_weeks)
df["BMI"] = pd.to_numeric(df["BMI_raw"], errors="coerce")
if "Age_raw" in df.columns: df["Age"] = pd.to_numeric(df["Age_raw"], errors="coerce")
df["w"]   = pd.to_numeric(df["w"], errors="coerce").fillna(1.0).clip(lower=0)

df = df[["id","Y","GA","BMI","w"] + (["Age"] if "Age" in df.columns else [])].dropna().copy()
print(f"样本量：{len(df)}；孕妇个体数：{df['id'].nunique()}；有效权重和={df['w'].sum():.1f}")

# SV + logit
df["Y_star"] = smithson_verkuilen(df["Y"])
df["logitY"] = np.log(df["Y_star"]/(1-df["Y_star"]))

# 加权中心化
GA_mu  = wmean(df["GA"],  df["w"])
BMI_mu = wmean(df["BMI"], df["w"])
df["GA_c"]  = df["GA"]  - GA_mu
df["BMI_c"] = df["BMI"] - BMI_mu
use_age = "Age" in df.columns
if use_age:
    Age_mu = wmean(df["Age"], df["w"])
    df["Age_c"] = df["Age"] - Age_mu

# 线性 LMM
X_lin = pd.DataFrame({"Intercept":1.0, "GA_c":df["GA_c"], "BMI_c":df["BMI_c"]})
if use_age: X_lin["Age_c"] = df["Age_c"]
res_lin = fit_mixedlm_weighted(df["logitY"].values, X_lin.values, df["id"], df["w"].values)
k_lin   = len(res_lin.params)
aic_lin = AIC(res_lin.llf, k_lin); bic_lin = BIC(res_lin.llf, k_lin, len(df))
print(f"[线性LMM·带权] logLik={res_lin.llf:.3f}  AIC={aic_lin:.2f}  BIC={bic_lin:.2f}")

# GAMM：cr样条+随机截距
best = None
df_grid = [4,5,6]
for df_ga in df_grid:
    for df_b in df_grid:
        formula = f"1 + cr(GA_c, df={df_ga}) + cr(BMI_c, df={df_b})" + (" + Age_c" if use_age else "")
        try:
            X_try = dmatrix(formula, df, return_type="dataframe")
            res   = fit_mixedlm_weighted(df["logitY"].values, X_try.values, df["id"], df["w"].values)
            k     = len(res.params)
            aic   = AIC(res.llf, k); bic = BIC(res.llf, k, len(df))
            print(f"  [GAMM·带权] df_GA={df_ga}, df_BMI={df_b}  logLik={res.llf:.3f}  AIC={aic:.2f}")
            if (best is None) or (aic < best["aic"]):
                best = dict(res=res, aic=aic, bic=bic, k=k, df_ga=df_ga, df_bmi=df_b,
                            formula=formula, design_info=X_try.design_info)
        except Exception as e:
            print(f"  [跳过] df_GA={df_ga}, df_BMI={df_b} 收敛失败：{e}")

# 退路：孕周样条+BMI 线性
fallback = False
if best is None:
    fallback = True
    for df_ga in df_grid:
        formula = f"1 + cr(GA_c, df={df_ga}) + BMI_c" + (" + Age_c" if use_age else "")
        try:
            X_try = dmatrix(formula, df, return_type="dataframe")
            res   = fit_mixedlm_weighted(df["logitY"].values, X_try.values, df["id"], df["w"].values)
            k     = len(res.params)
            aic   = AIC(res.llf, k); bic = BIC(res.llf, k, len(df))
            print(f"  [GAMM-lite·带权] df_GA={df_ga}  logLik={res.llf:.3f}  AIC={aic:.2f}")
            if (best is None) or (aic < best["aic"]):
                best = dict(res=res, aic=aic, bic=bic, k=k, df_ga=df_ga, df_bmi=None,
                            formula=formula, design_info=X_try.design_info)
        except Exception as e:
            print(f"  [跳过] GAMM-lite df_GA={df_ga} 收敛失败：{e}")

if best is None:
    raise SystemExit("所有 GAMM 设定均收敛失败，请检查数据或降低样条自由度。")

res_gamm = best["res"]; design_info = best["design_info"]
print("\n[GAMM 最优·带权] 公式:", best["formula"])
print(f"[GAMM 最优·带权] logLik={res_gamm.llf:.3f}  AIC={best['aic']:.2f}")

# LRT
df_diff = best["k"] - k_lin
LR, dfd, pval = lrt(res_lin, res_gamm, df_diff=df_diff)
print(f"\n[非线性整体检验·带权] 线性LMM → GAMM：LR={LR:.3f}, df={dfd}, p={pval:.4g}")

#固定效应表
X_train = dmatrix(best["formula"], df, return_type="dataframe")
p_fix   = X_train.shape[1]
coef = pd.Series(res_gamm.params[:p_fix], index=[f"FE_{i}" for i in range(p_fix)], name="Estimate")
se   = pd.Series(res_gamm.bse[:p_fix],    index=coef.index,  name="SE")
z    = pd.Series(coef.values/se.values,   index=coef.index,  name="z")
p    = pd.Series(2*stats.norm.sf(np.abs(z.values)), index=coef.index, name="p")
ciL  = pd.Series(coef.values - stats.norm.ppf(0.975)*se.values, index=coef.index, name="CI_low")
ciH  = pd.Series(coef.values + stats.norm.ppf(0.975)*se.values, index=coef.index, name="CI_high")
pd.concat([coef,se,z,p,ciL,ciH], axis=1).to_excel(out_dir/"GAMM_coef_table_weighted.xlsx")

with open(out_dir/"GAMM_model_summary_weighted.txt","w",encoding="utf-8") as f:
    f.write(
        f"数据源：{used_path}/{used_sheet}\n"
        f"样本量：{len(df)}；孕妇数：{df['id'].nunique()}；权重和：{df['w'].sum():.1f}\n"
        f"线性LMM·带权：logLik={res_lin.llf:.3f}, AIC={aic_lin:.2f}\n"
        f"GAMM最优·带权：{best['formula']}\n"
        f"logLik={res_gamm.llf:.3f}, AIC={best['aic']:.2f}\n"
        f"LRT(线性→GAMM)：LR={LR:.3f}, df={dfd}, p={pval:.4g}\n"
        f"说明：{'退路模型（孕周样条+BMI线性）' if fallback else '二维样条（孕周+BMI）'}\n"
    )

#部分效应图
def partial_plot(var_name, x_grid, GA_med, BMI_med, Age_med=None, title="", fn="plot.png"):
    new = pd.DataFrame({
        "GA_c":  np.repeat(GA_med, len(x_grid)),
        "BMI_c": np.repeat(BMI_med, len(x_grid))
    })
    new[var_name] = x_grid
    if use_age: new["Age_c"] = Age_med
    X_new = build_design_matrices([design_info], new)[0]
    X_new = np.asarray(X_new)
    beta = res_gamm.params[:p_fix]
    cov  = res_gamm.cov_params()[:p_fix, :p_fix]
    lin      = X_new @ beta
    var_lin  = np.einsum("ij,jk,ik->i", X_new, cov, X_new)
    se_lin   = np.sqrt(np.maximum(var_lin, 0))
    zcrit    = stats.norm.ppf(0.975)
    lo_prob  = inv_logit(lin - zcrit*se_lin)
    hi_prob  = inv_logit(lin + zcrit*se_lin)
    y_prob   = inv_logit(lin)
    plt.figure(figsize=(7,5))
    if var_name == "GA_c":
        x_axis = x_grid + wmean(df["GA"], df["w"])
        plt.xlabel("检测孕周（周）")
    else:
        x_axis = x_grid + wmean(df["BMI"], df["w"])
        plt.xlabel("BMI")
    plt.plot(x_axis, y_prob, label="预测均值")
    plt.fill_between(x_axis, lo_prob, hi_prob, alpha=0.25, label="95%CI")
    plt.ylabel("预测Y浓度（比例）")
    plt.title(title)
    plt.legend(); plt.tight_layout(); plt.savefig(out_dir/fn, dpi=220); plt.close()

# 用带权“中位数”和范围
def wmedian(x, w):
    x = pd.to_numeric(x, errors="coerce"); w = pd.to_numeric(w, errors="coerce").clip(lower=0)
    idx = np.argsort(x.values); xv, wv = x.values[idx], w.values[idx]
    cw = np.cumsum(wv); total = cw[-1] if len(cw)>0 else np.nan
    if not np.isfinite(total) or total<=0: return np.nan
    return xv[np.searchsorted(cw, 0.5*total)]

GA_med  = wmedian(df["GA_c"], df["w"]); BMI_med = wmedian(df["BMI_c"], df["w"]); Age_med = wmedian(df["Age_c"], df["w"]) if use_age else None
GA_lo, GA_hi   = df["GA_c"].quantile(0.02),  df["GA_c"].quantile(0.98)
BMI_lo, BMI_hi = df["BMI_c"].quantile(0.02), df["BMI_c"].quantile(0.98)
GA_grid  = np.linspace(GA_lo,  GA_hi,  150)
BMI_grid = np.linspace(BMI_lo, BMI_hi, 150)

partial_plot("GA_c",  GA_grid,  GA_med, BMI_med, Age_med,
             title="GAMM（带权/兼容）部分效应：孕周", fn="GAMM_w_partial_GA.png")
partial_plot("BMI_c", BMI_grid, GA_med,  BMI_med, Age_med,
             title="GAMM（带权/兼容）部分效应：BMI", fn="GAMM_w_partial_BMI.png")

print("\n[GAMM·带权] 输出：", out_dir.resolve())


\end{lstlisting}
\end{appendices}

 \noindent \textbf{0.4}这是使用线性混合模型（LMM）分析同一孕期数据的核心程序。它计算加权均值，并使用混合效应方法拟合变换后的数据（logit转换）。输出包括模型评估所需的AIC、BIC和对数似然值。用于与GAMM模型对比。
\begin{appendices}
\begin{lstlisting}

import re
import numpy as np
import pandas as pd
import statsmodels.api as sm
import statsmodels.formula.api as smf
from patsy import dmatrices
from scipy import stats
import matplotlib, matplotlib.pyplot as plt
from pathlib import Path
import warnings


matplotlib.rcParams["font.sans-serif"] = ["Microsoft YaHei", "SimHei", "Arial Unicode MS", "DejaVu Sans"]
matplotlib.rcParams["axes.unicode_minus"] = False

CANDIDATES = [
    ("附件-男胎检测数据-处理.xlsx", "附件-男胎检测数据-处理"),
    ("附件.xlsx", "男胎检测数据"),
]
COL_ID, COL_Y, COL_GA, COL_BMI, COL_AGE, COL_W = "孕妇代码", "Y染色体浓度", "检测孕周", "孕妇BMI", "年龄", "w_Q"
out_dir = Path("./model_outputs_lmm_w"); out_dir.mkdir(parents=True, exist_ok=True)

def parse_gestation_to_weeks(x):
    if pd.isna(x): return np.nan
    s = str(x).strip()
    try: return float(s)
    except: pass
    m = re.search(r'(\d+)\s*(?:周|w|W)?\s*\+?\s*(\d+)?\s*(?:天|d|D)?', s)
    if m:
        w = float(m.group(1)); d = float(m.group(2)) if m.group(2) else 0.0
        return w + d/7.0
    m2 = re.search(r'(\d+)', s)
    return float(m2.group(1)) if m2 else np.nan

def smithson_verkuilen(y):
    y = y.clip(0, 1); n = y.notna().sum()
    return (y*(n-1) + 0.5) / n

def wmean(x, w):
    x = pd.to_numeric(x, errors="coerce"); w = pd.to_numeric(w, errors="coerce").clip(lower=0)
    return np.nansum(w*x)/np.nansum(w) if np.nansum(w)>0 else np.nan

def mixedlm_lrt(res_small, res_large, df_diff=None):
    LR = 2.0*(res_large.llf - res_small.llf)
    if df_diff is None:
        df_diff = len(res_large.params) - len(res_small.params)
    p = stats.chi2.sf(LR, df_diff)
    return LR, df_diff, p

# 读取
last_err = None
for path, sheet in CANDIDATES:
    try:
        df = pd.read_excel(path, sheet_name=sheet)
        used_path, used_sheet = path, sheet
        break
    except Exception as e:
        last_err = e
else:
    raise SystemExit(f"读取失败：{last_err}")

for c in [COL_ID, COL_Y, COL_GA, COL_BMI]:
    if c not in df.columns:
        raise SystemExit(f"找不到列：{c}（数据源：{used_path}/{used_sheet}）")
if COL_W not in df.columns:
    print("[提示] 未找到 w_Q 列，默认权重=1")
    df[COL_W] = 1.0

use_cols = [COL_ID, COL_Y, COL_GA, COL_BMI, COL_W] + ([COL_AGE] if COL_AGE in df.columns else [])
df = df[use_cols].copy()
df = df.rename(columns={COL_ID:"id", COL_Y:"Y", COL_GA:"GA_raw", COL_BMI:"BMI_raw", COL_W:"w"})
if COL_AGE in df.columns: df = df.rename(columns={COL_AGE:"Age_raw"})

df["GA"]  = df["GA_raw"].apply(parse_gestation_to_weeks)
df["BMI"] = pd.to_numeric(df["BMI_raw"], errors="coerce")
if "Age_raw" in df.columns: df["Age"] = pd.to_numeric(df["Age_raw"], errors="coerce")
df["w"]   = pd.to_numeric(df["w"], errors="coerce").fillna(1.0).clip(lower=0)

keep = ["id","Y","GA","BMI","w"] + (["Age"] if "Age" in df.columns else [])
df = df[keep].dropna().copy()

# SV + logit
df["Y_star"] = smithson_verkuilen(df["Y"])
df["logitY"] = np.log(df["Y_star"]/(1-df["Y_star"]))

# 加权中心化
GA_mu  = wmean(df["GA"],  df["w"])
BMI_mu = wmean(df["BMI"], df["w"])
df["GA_c"]  = df["GA"]  - GA_mu
df["BMI_c"] = df["BMI"] - BMI_mu
use_age = "Age" in df.columns
if use_age:
    Age_mu = wmean(df["Age"], df["w"])
    df["Age_c"] = df["Age"] - Age_mu

print(f"样本量：{len(df)}；孕妇个体数：{df['id'].nunique()}；有效权重和={df['w'].sum():.1f}")

#带权拟合
def fit_mlm(formula_str):
    w = np.asarray(df["w"].values, dtype=float)
    try:
        md = smf.mixedlm(formula_str, data=df, groups=df["id"], re_formula="1", weights=w)
        res = md.fit(reml=False, method="lbfgs")
        # 保持统一接口
        try:
            res._exog_names = list(res.model.exog_names)
        except Exception:
            pass
        return res
    except (TypeError, ValueError):
        print(f"[警告] MixedLM 不支持 weights，公式 {formula_str} 改用 √w 伪加权。")
        # 用 patsy 生成设计矩阵（含截距，并能拿到列名）
        y_mat, X_mat = dmatrices(formula_str, df, return_type="dataframe")
        y = np.asarray(y_mat).ravel()
        X = np.asarray(X_mat); exog_names = list(X_mat.columns)
        sw = np.sqrt(np.clip(w, 1e-8, None))
        y_t = y * sw
        X_t = X * sw[:, None]
        Z_t = np.ones((len(df), 1)) * sw[:, None]   # 随机截距也乘 √w
        md = sm.MixedLM(y_t, X_t, groups=df["id"], exog_re=Z_t)
        res = md.fit(reml=False, method="lbfgs")
        # 把名字挂到结果上
        res._exog_names = exog_names
        return res

fit0 = fit_mlm("logitY ~ 1")
fit1 = fit_mlm("logitY ~ 1 + GA_c")
fit2 = fit_mlm("logitY ~ 1 + GA_c + BMI_c")
if use_age:
    fit3 = fit_mlm("logitY ~ 1 + GA_c + BMI_c + Age_c")
    final_fit, final_name = fit3, "M3: GA + BMI + Age（带权/兼容）"
else:
    final_fit, final_name = fit2, "M2: GA + BMI（带权/兼容）"

print("\n=== 最终模型摘要（{}） ===".format(final_name))
print(final_fit.summary())

# LRT
LR_10 = mixedlm_lrt(fit0, fit1, df_diff=1)
LR_21 = mixedlm_lrt(fit1, fit2, df_diff=1)
print("\n=== 似然比检验（带权/兼容） ===")
print(f"M0→M1（+孕周）: LR={LR_10[0]:.3f}, df={LR_10[1]}, p={LR_10[2]:.4g}")
print(f"M1→M2（+BMI） : LR={LR_21[0]:.3f}, df={LR_21[1]}, p={LR_21[2]:.4g}")
if use_age:
    LR_32 = mixedlm_lrt(fit2, fit3, df_diff=1)
    print(f"M2→M3（+年龄）: LR={LR_32[0]:.3f}, df={LR_32[1]}, p={LR_32[2]:.4g}")

# 边际/条件R²（简化近似）
try:
    var_random = float(getattr(final_fit, "cov_re", np.array([[np.nan]])).iloc[0,0])
except Exception:
    var_random = float(np.asarray(getattr(final_fit, "cov_re", np.array([[np.nan]])))[0,0])
var_resid = float(final_fit.scale)
var_fitted = float(np.var(final_fit.fittedvalues, ddof=1))
denom = var_fitted + var_resid if np.isfinite(var_fitted + var_resid) else np.nan
r2_c = var_fitted/denom if denom and denom>0 else np.nan
r2_m = max(0.0, r2_c - var_random/(var_random + var_resid + 1e-9)) if np.isfinite(r2_c) else np.nan
print("\n=== R²（Nakagawa & Schielzeth，带权/兼容） ===")
print(f"边际R²：{r2_m:.4f}；条件R²：{r2_c:.4f}")
print("方差组成（近似）：", {'var_random': var_random, 'var_resid': var_resid})

# 等效孕周 & OR 
# 固定效应的名字与向量
fe_names = getattr(final_fit, "_exog_names", None)
if fe_names is None and hasattr(final_fit, "model") and hasattr(final_fit.model, "exog_names"):
    fe_names = list(final_fit.model.exog_names)

k_fe = int(getattr(final_fit, "k_fe", len(fe_names) if fe_names is not None else 0))
beta_all = np.asarray(final_fit.params)
beta_fe  = beta_all[:k_fe]                      
cov_all  = np.asarray(final_fit.cov_params())
cov_fe   = cov_all[:k_fe, :k_fe]                 

if fe_names is not None and "GA_c" in fe_names and "BMI_c" in fe_names:
    idx_ga  = fe_names.index("GA_c")
    idx_bmi = fe_names.index("BMI_c")
    b_ga  = float(beta_fe[idx_ga])
    b_bmi = float(beta_fe[idx_bmi])
    cov2  = cov_fe[[idx_ga, idx_bmi]][:, [idx_ga, idx_bmi]]

    theta = -b_bmi / b_ga
    d_b1 = b_bmi / (b_ga**2); d_b2 = -1.0 / b_ga
    grad = np.array([[d_b1, d_b2]])
    var_theta = float(np.squeeze(grad @ cov2 @ grad.T))
    se = np.sqrt(max(var_theta, 0))
    z = stats.norm.ppf(0.975); lo, hi = theta - z*se, theta + z*se

    print("\n=== ‘等效孕周’（带权/兼容） ===")
    print(f"点估计：{theta:.3f} 周；95%CI [{lo:.3f}, {hi:.3f}]；SE={se:.3f}")
    print("\n=== 胜算比 OR（logit 尺度） ===")
    print(f"孕周 +1 周：OR={np.exp(b_ga):.3f}")
    print(f"BMI  +1   ：OR={np.exp(b_bmi):.3f}")
else:
    print("\n[提示] 无法定位 GA_c/BMI_c 的系数（因列名丢失），跳过‘等效孕周/OR’。")

#  诊断与部分效应 
def inv_logit(z): return 1.0/(1.0+np.exp(-z))
plt.figure(figsize=(7,5))
plt.scatter(final_fit.fittedvalues, final_fit.resid, s=10, alpha=0.6)
plt.axhline(0, color="k", lw=1); plt.xlabel("拟合值（logitY）"); plt.ylabel("残差")
plt.title("残差-拟合图（LMM·带权/兼容）")
plt.tight_layout(); plt.savefig(out_dir/"resid_vs_fitted_w.png", dpi=220)

# 用 names 来构造预测矩阵
def predict_curve(var_name, grid):
    # 仅用固定效应名/参数预测
    if fe_names is None:
        return None, None
    base = {nm:0.0 for nm in fe_names}
    if "Intercept" in base: base["Intercept"] = 1.0

    X_rows = []
    for val in grid:
        row = base.copy()
        if var_name in row:
            row[var_name] = val
        X_rows.append([row[nm] for nm in fe_names])
    X = np.asarray(X_rows)

    yhat = inv_logit(X @ beta_fe)   # 用固定效应参数
    # 反中心化回原尺度
    if var_name == "GA_c":
        x_axis = grid + GA_mu
    elif var_name == "BMI_c":
        x_axis = grid + BMI_mu
    else:
        x_axis = grid
    return x_axis, yhat


GA_grid  = np.linspace(df["GA_c"].quantile(0.02),  df["GA_c"].quantile(0.98), 120)
BMI_grid = np.linspace(df["BMI_c"].quantile(0.02), df["BMI_c"].quantile(0.98), 120)

xg, yg = predict_curve("GA_c", GA_grid)
if xg is not None:
    plt.figure(figsize=(7,5))
    plt.plot(xg, yg); plt.xlabel("检测孕周（周）"); plt.ylabel("预测Y浓度（比例）")
    plt.title("部分效应：孕周（带权/兼容）")
    plt.tight_layout(); plt.savefig(out_dir/"partial_effect_GA_w.png", dpi=220)

xb, yb = predict_curve("BMI_c", BMI_grid)
if xb is not None:
    plt.figure(figsize=(7,5))
    plt.plot(xb, yb); plt.xlabel("BMI"); plt.ylabel("预测Y浓度（比例）")
    plt.title("部分效应：BMI（带权/兼容）")
    plt.tight_layout(); plt.savefig(out_dir/"partial_effect_BMI_w.png", dpi=220)
plt.close("all")

if fe_names is None:
    fe_names = [f"x{i+1}" for i in range(k_fe)]
fe_se = np.asarray(final_fit.bse)[:k_fe]

coef = pd.DataFrame({"name": fe_names, "Estimate": beta_fe})
z    = coef["Estimate"].values / fe_se
p    = 2*stats.norm.sf(np.abs(z))
ciL  = coef["Estimate"].values - stats.norm.ppf(0.975)*fe_se
ciH  = coef["Estimate"].values + stats.norm.ppf(0.975)*fe_se

out_coef = pd.DataFrame({
    "Estimate": coef["Estimate"].values,
    "SE": fe_se,
    "z": z,
    "p": p,
    "CI_low": ciL,
    "CI_high": ciH
}, index=fe_names)
out_coef.to_excel(out_dir/"LMM_coef_table_weighted_compat.xlsx")


with open(out_dir/"model_summary_weighted_compat.txt","w",encoding="utf-8") as f:
    f.write(
        f"数据源：{used_path}/{used_sheet}\n"
        f"样本量：{len(df)}；孕妇数：{df['id'].nunique()}；权重和：{df['w'].sum():.1f}\n"
        f"最终模型：{final_name}\n"
        f"边际R²：{r2_m:.4f}；条件R²：{r2_c:.4f}\n"
        f"随机截距方差：{var_random:.6f}；残差方差：{var_resid:.6f}\n"
    )

print("\n[LMM·带权/兼容] 输出：", out_dir.resolve())


\end{lstlisting}
\end{appendices}

 \noindent \textbf{0.5}这是解决问题二所用核心代码，进行数据预处理、阈值计算和对各种基因标记进行软质量控制标记，处理与Y染色体浓度及其他测量相关的数据，应用质量控制的阈值。
\begin{appendices}
\begin{lstlisting}

###问题2_核心代码


#单人 Ti 构造
def construct_Ti(df, res, design_info, work_for_centering):
    GA_mean  = work_for_centering["GA"].mean()
    BMI_mean = work_for_centering["BMI"].mean()
    Age_mean = work_for_centering["Age"].mean() if "Age" in work_for_centering.columns else None

    rows = []
    for pid, g in df.sort_values([COL_ID, COL_GA]).groupby(COL_ID):
        g = g.sort_values(COL_GA)
        ga = g[COL_GA].to_numpy(float)
        yy = g[COL_Y].to_numpy(float)
        bmi0 = float(g[COL_BMI].iloc[0])
        age0 = float(g[COL_AGE].iloc[0]) if (COL_AGE in g.columns) else None

        thr = 0.04
        cross_T = None
        for j in range(1, len(ga)):
            if (yy[j-1] < thr) and (yy[j] >= thr):
                t1, t2 = ga[j-1], ga[j]
                y1, y2 = yy[j-1], yy[j]
                w = 0.5 if (y2==y1) else np.clip((thr - y1)/(y2 - y1), 0.0, 1.0)
                cross_T = t1 + w*(t2 - t1)
                break

        if cross_T is not None:
            rows.append(dict(id=pid, BMI0=bmi0, Age0=age0, T=float(cross_T),
                             source="interp", n_tests=len(g), first_ga=float(ga[0]), 
                             last_ga=float(ga[-1])))
            continue

        #无跨越 → 用带权 GAMM 的固定效应解根
        def f(t):
            lin = predict_logit_mean(design_info, res,
                                     GA=np.array([t]),
                                     BMI=np.array([bmi0]),
                                     Age=np.array([age0]) if (age0 is not None) else None,
                                     GA_mean=GA_mean, BMI_mean=BMI_mean, Age_mean=Age_mean)
            return inv_logit(lin) - thr

        t_lo, t_hi = 8.0, 32.0
        f_lo, f_hi = f(t_lo), f(t_hi)
        if f_lo >= 0:
            Ti = max(10.0, float(ga[0])) if len(ga)>0 else 10.0
        elif f_hi <= 0:
            Ti = max(32.0, float(ga[-1])) if len(ga)>0 else 32.0
        else:
            Ti = float(brentq(f, t_lo, t_hi, maxiter=200))
        rows.append(dict(id=pid, BMI0=bmi0, Age0=age0, T=Ti,
                         source="gamm", n_tests=len(g), first_ga=float(ga[0]), last_ga=float(ga[-1])))
    return pd.DataFrame(rows).sort_values("BMI0").reset_index(drop=True)


#DP 分段
def precompute_sums_logT(Ti_df, weights=None):
    z = np.log(np.clip(Ti_df["T"].values, 1e-6, None))
    if weights is None:
        w = np.ones_like(z)
    else:
        w = np.asarray(weights, float)
    W   = np.concatenate([[0.0], np.cumsum(w)])
    S   = np.concatenate([[0.0], np.cumsum(w * z)])
    S2  = np.concatenate([[0.0], np.cumsum(w * z * z)])
    return z, w, W, S, S2


#加权核ECDF
ALPHA_KERNEL = ALPHA_KERNEL
def kernel_ecdf_factory_weighted(samples, weights=None, alpha=ALPHA_KERNEL):
    xs = np.asarray(samples, float)
    n  = len(xs)
    if n == 0: return lambda t: 0.0
    if weights is None:
        w = np.ones(n, float)
    else:
        w = np.asarray(weights, float).clip(min=0.0)
        if w.sum() <= 0: w = np.ones(n, float)
    mu = np.sum(w*xs)/np.sum(w)
    var = np.sum(w*(xs-mu)**2)/np.sum(w)
    sd = np.sqrt(max(var, 1e-12))
    h_silverman = 1.06 * sd * (n ** (-1/5))
    h = max(H_MIN, alpha * h_silverman)
    w_norm = w / np.sum(w); inv_h = 1.0 / h
    def Fhat(t):
        t_arr = np.atleast_1d(t).astype(float)
        Z = (t_arr[:, None] - xs[None, :]) * inv_h
        vals = stats.norm.cdf(Z) @ w_norm
        return float(vals[0]) if np.isscalar(t) else vals
    return Fhat

#风险函数
def expected_risk_enhanced(F, t0, delta=DELTA, group_stats=None, 
                           group_idx=None, total_groups=None):
    t_max = max(40.0, 28.0 + 6.0)
    risk = LAMBDA_FAIL * (1.0 - F(t0))
    if group_stats and group_idx is not None and total_groups is not None:
        mean_T = group_stats['mean_T']
        std_T = group_stats['std_T']
        mean_BMI = group_stats['mean_BMI']
        if mean_BMI > 35:      
            early_penalty, late_penalty = 0.5, 2.0
        elif mean_BMI < 28:    
            early_penalty, late_penalty = 1.5, 0.8
        else:                  
            early_penalty, late_penalty = 1.0, 1.0
        risk *= early_penalty
        variability_factor = 1.0 + 0.1 * (std_T / max(mean_T, 1))
        t_prev = 0.0; tk = t0
        while tk <= t_max:
            mass = max(F(tk) - F(t_prev), 0.0)
            base_weight = stage_weight(tk)
            if tk > 20: base_weight *= late_penalty
            risk += base_weight * variability_factor * mass
            t_prev = tk; tk += delta
        tail = max(1.0 - F(t_prev), 0.0)
        if tail > 0:
            risk += stage_weight(t_max + 1e-6) * late_penalty * variability_factor * tail
    else:
        t_prev = 0.0; tk = t0
        while tk <= t_max:
            mass = max(F(tk) - F(t_prev), 0.0)
            risk += stage_weight(tk) * mass
            t_prev = tk; tk += delta
        tail = max(1.0 - F(t_prev), 0.0)
        if tail > 0:
            risk += stage_weight(t_max + 1e-6) * tail
    return float(risk)


#最佳时点
def optimize_t_star_enhanced(F, group_stats, group_idx, total_groups):
    mean_T = group_stats['mean_T']; std_T = group_stats['std_T']
    median_T = group_stats['median_T']; mean_BMI = group_stats['mean_BMI']
    if mean_BMI < 28:
        win_center = max(WIN_MIN, min(median_T - std_T, 14))
        win_width = min(6, 1.5 * std_T)
    elif mean_BMI > 35:
        win_center = min(WIN_MAX - 2, max(median_T, 16))
        win_width = min(8, 2 * std_T)
    else:
        win_center = max(WIN_MIN + 2, min(median_T - 0.3 * std_T, WIN_MAX - 3))
        win_width = min(7, 1.8 * std_T)
    search_min = max(WIN_MIN, win_center - win_width/2)
    search_max = min(WIN_MAX, win_center + win_width/2)
    if search_max - search_min < 4:
        if mean_BMI < 30: search_max = min(WIN_MAX, search_min + 4)
        else:             search_min = max(WIN_MIN, search_max - 4)
    t_grid = np.arange(search_min, search_max + 1e-9, T_GRID_STEP)
    risks = [expected_risk_enhanced(F, t, group_stats=group_stats,
        group_idx=group_idx, total_groups=total_groups) for t in t_grid]
    if len(risks) == 0: return float(mean_T), 1.0
    k = int(np.argmin(risks))
    return float(t_grid[k]), float(risks[k])

\end{lstlisting}
\end{appendices}

 \noindent \textbf{0.6}这是解决问题三所用核心代码，继续运用了问题一的广义加性模型和正交化技术分析孕期相关数据，处理BMI、年龄、身高和体重等特征，预测Y染色体浓度。
\begin{appendices}
\begin{lstlisting}

###问题3_核心代码



#拟合 GAMM
def fit_gamm(df):
    def sv(y):
        y = np.clip(y, 0, 1); n = np.isfinite(y).sum()
        return (y*(n-1)+0.5)/max(n,1)

    work = df.copy()
    work["GA"]  = work[COL_GA]
    work["BMI"] = work[COL_BMI]
    if COL_AGE in work.columns: work["Age"] = work[COL_AGE]
    if COL_HT  in work.columns: work["Ht"]  = work[COL_HT]
    if COL_WT  in work.columns: work["Wt"]  = work[COL_WT]

    work, beta_hw = hw_orthogonalize(work) 

    work["Y_star"] = sv(work[COL_Y].values)
    work["logitY"] = np.log(work["Y_star"]/(1.0 - work["Y_star"]))
    work["GA_c"]   = work["GA"]  - work["GA"].mean()
    work["BMI_c"]  = work["BMI"] - work["BMI"].mean()
    work["rHT_c"]  = work["rHT"] - work["rHT"].mean()
    work["rWT_c"]  = work["rWT"] - work["rWT"].mean()
    if "Age" in work.columns: work["Age_c"] = work["Age"] - work["Age"].mean()

    formula = f"1 + cr(GA_c, df={DF_GA}) * cr(BMI_c, df={DF_BMI})" \
              f" + cr(rHT_c, df={DF_RHT}) + cr(rWT_c, df={DF_RWT})"
    if "Age_c" in work.columns:
        formula += f" + cr(Age_c, df={DF_AGE})"

    try:
        X = dmatrix(formula, work, return_type="dataframe")
        md = sm.MixedLM(work["logitY"].values, X.values,
                        groups=work[COL_ID].astype(str).values,
                        exog_re=np.ones((len(work),1)))
        res = md.fit(reml=False, method="lbfgs")
    except Exception as e:
        raise RuntimeError(f"GAMM 拟合失败：{e}\n公式：{formula}")

    resid_var = float(res.scale)
    try:
        group_var = float(res.cov_re.iloc[0,0]) if hasattr(res,"cov_re") else float(res.cov_re[0,0])
    except Exception:
        group_var = float(np.squeeze(np.array(res.cov_re)))

    print(f"[GAMM] 公式：{formula}")
    print(f"[GAMM] 残差方差={resid_var:.4f}，组方差={group_var:.4f}")
    return res, X.design_info, work, resid_var, group_var, beta_hw


#构造个体F_i(t)
def build_F_from_gamm(sub, res, design_info, work, t_grid, beta_hw):
    GA_mean  = work["GA"].mean()
    BMI_mean = work["BMI"].mean()
    Age_mean = work["Age"].mean() if "Age" in work.columns else None
    rHT_mean = work["rHT"].mean()
    rWT_mean = work["rWT"].mean()
    thr = logit(Y_THRESH)

    N = len(sub); T = len(t_grid)
    F_mat = np.zeros((N, T), dtype=float)

    for idx, row in enumerate(sub.itertuples(index=False)):
        bmi0 = getattr(row, "BMI0")
        age0 = getattr(row, COL_AGE) if (COL_AGE in sub.columns) else None
        ht0  = getattr(row, COL_HT)  if (COL_HT  in sub.columns) else None
        wt0  = getattr(row, COL_WT)  if (COL_WT  in sub.columns) else None

        rht_vec, rwt_vec = hw_orth_apply(bmi0, ht0, wt0, beta_hw, len(t_grid))

        new = {
            "GA_c":  (t_grid - GA_mean),
            "BMI_c": np.full_like(t_grid, bmi0 - BMI_mean),
            "rHT_c": (rht_vec - rHT_mean),
            "rWT_c": (rwt_vec - rWT_mean)
        }
        if "Age" in work.columns:
            new["Age_c"] = np.full_like(t_grid, (age0 if age0 is not None else Age_mean) - Age_mean)

        mu = predict_mu(design_info, res, new)
        p  = inv_logit(mu - thr)
        F  = np.maximum.accumulate(np.clip(p, 0.0, 1.0))
        F_mat[idx,:] = F
    return F_mat


#DP（风险段成本）
def build_time_grid():
    t0 = 0.0
    t_grid = np.arange(t0, T_MAX + 1e-9, DT)
    t_candidates = np.arange(WIN_MIN, WIN_MAX + 1e-9, DT)
    idx = lambda t: int(round((t - t0)/DT))
    tk_index = {}
    for t in t_candidates:
        items = []
        tk = t; prev = 0.0
        while tk <= T_MAX + 1e-9:
            items.append((idx(tk), idx(prev), stage_weight(tk)))
            prev = tk; tk += DELTA
        tk_index[t] = items
    return t_grid, t_candidates, tk_index


def dp_partition_risk(sub_sorted, F_sorted, t_grid, t_candidates, tk_index,
                      gmax=G_MAX, nmin=N_MIN_GROUP, gamma=GAMMA_STRUCT):
    N, T = F_sorted.shape
    prefix = np.vstack([np.zeros((1,T)), np.cumsum(F_sorted, axis=0)])
    def seg_sum(i,j,ti): return prefix[j+1,ti] - prefix[i,ti]

    INF = 1e100
    cost = np.full((N,N), np.inf, float)
    tstar_idx = np.full((N,N), -1, int)

    for i in range(N):
        for j in range(i,N):
            if (j-i+1) < nmin: continue
            best_val = INF; best_idx = -1
            n_g = (j-i+1)
            for t in t_candidates:
                acc = 0.0
                for (it, ip, wt) in tk_index[t]:
                    acc += wt * (seg_sum(i,j,it) - seg_sum(i,j,ip))
                it0 = int(round((t-0.0)/DT))
                acc += (1.0 - seg_sum(i,j,it0)/max(n_g,1)) * n_g * LAMBDA_FAIL
                if acc < best_val:
                    best_val = acc; best_idx = it0
            cost[i,j] = best_val; tstar_idx[i,j] = best_idx

    DP = np.full((gmax+1, N), INF, float); PREV = np.full((gmax+1, N), -1, int)
    for j in range(N): DP[1,j] = cost[0,j]; PREV[1,j] = -1
    for k in range(2,gmax+1):
        for j in range(k-1,N):
            best=INF; best_i=-1
            for i in range(k-1,j+1):
                v = DP[k-1,i-1] + cost[i,j]
                if v < best: best=v; best_i=i
            DP[k,j]=best; PREV[k,j]=best_i

    logN = np.log(N)
    crit_list=[]
    for k in range(1,gmax+1):
        raw = DP[k,N-1]; pen = gamma*(k-1)*logN
        crit_list.append((k, raw, raw+pen))
        print(f"  k={k}: 原始成本={raw:.2f}, 惩罚={pen:.2f}, 总准则={raw+pen:.2f}")
    best_k, _, _ = min(crit_list, key=lambda x: x[2])

    segs=[]; j=N-1; k=best_k
    while k>=1:
        i=PREV[k,j]; i=0 if i==-1 else i
        segs.append((i,j)); j=i-1; k-=1
    segs=list(reversed(segs))
    return best_k, segs, crit_list, cost, tstar_idx


def summarize_groups(sub_sorted, F_sorted, segs, t_grid, tstar_idx, tk_index):
    rows=[]
    for g,(s,e) in enumerate(segs, start=1):
        it=tstar_idx[s,e]; t_star=t_grid[it]
        Fg = F_sorted[s:e+1,:]
        one_shot=float(Fg[:,it].mean()); retest=1.0-one_shot
        A=Fg.sum(axis=0); acc=0.0
        for (itk, ip, wt) in tk_index[t_star]:
            acc += wt*(A[itk]-A[ip])
        acc += (1.0 - A[it]/max((e-s+1),1))*(e-s+1)*LAMBDA_FAIL
        rows.append({
            "组别":g, "人数":(e-s+1),
            "BMI下界":float(sub_sorted["BMI0"].iloc[s]),
            "BMI上界":float(sub_sorted["BMI0"].iloc[e]),
            "最佳检测时点(周)":float(t_star),
            "一次达标率":one_shot, "重测率":retest, "风险值":float(acc)
        })
    return pd.DataFrame(rows)

\end{lstlisting}
\end{appendices}

 \noindent \textbf{0.7}这是问题四所用核心代码，特征提取并输入到所用的两个机器学习模型中进行训练。通过重采样方法解决类别不平衡的问题。
\begin{appendices}
\begin{lstlisting}

###问题4_核心代码



# 机器学习相关
from sklearn.model_selection import (
    train_test_split, cross_val_score, GridSearchCV,
    StratifiedKFold, cross_validate, GroupKFold
)
from sklearn.ensemble import (
    GradientBoostingClassifier, AdaBoostClassifier,
    RandomForestClassifier, ExtraTreesClassifier
)
from sklearn.metrics import (
    roc_auc_score, average_precision_score, roc_curve,
    precision_recall_curve, accuracy_score, precision_score,
    recall_score, f1_score, confusion_matrix, classification_report,
    brier_score_loss, log_loss, matthews_corrcoef, cohen_kappa_score
)
from sklearn.preprocessing import StandardScaler
from sklearn.impute import SimpleImputer

# 重采样
from imblearn.over_sampling import (
    SMOTE, BorderlineSMOTE, SVMSMOTE, RandomOverSampler, ADASYN
)
from imblearn.under_sampling import RandomUnderSampler
from imblearn.combine import SMOTEENN, SMOTETomek

import joblib
warnings.filterwarnings('ignore')

# 设置随机种子
RANDOM_STATE = 42
np.random.seed(RANDOM_STATE)

_MODEL_ZH = {
    "GradientBoosting": "梯度提升(GBDT)",
    "AdaBoost": "AdaBoost"
}
_DATASET_ZH = {
    "original": "原始数据",
    "random_oversample": "随机过采样",
    "random_undersample": "随机欠采样",
    "smote": "SMOTE",
    "borderline_smote": "边界SMOTE",
    "smoteenn": "SMOTEENN",
}
_METRIC_ZH = {
    "AUC": "AUC",
    "F1": "F1",
    "Recall": "召回率",
    "Precision": "精确率",
    "MCC": "MCC",
    "Specificity": "特异性",
}
def _zh_model(name): return _MODEL_ZH.get(name, name)
def _zh_dataset(name): return _DATASET_ZH.get(name, name)
def _zh_metric(name): return _METRIC_ZH.get(name, name)


class AneuploidyDetector:
    def __init__(self, config=None):
        self.config = config or self._get_default_config()
        self.models = {}
        self.results = {}
        self.feature_names = None
        self.sample_weights = None
        self.pregnant_groups = None
        self._setup_directories()

    def _get_default_config(self):
        return {
            'input_file': './附件-女胎处理_01处理.xlsx',
            'sheet_name': '附件-女胎处理',
            'abnormal_col': '染色体的非整倍体',
            'pregnant_id_col': '孕妇代码',  # 孕妇ID列
            'weight_col': 'w_Q',          # 权重列
            'output_dir': Path('./gbdt_enhanced_results'),
            'test_ratio': 0.3,
            'cv_folds': 5,
            'random_state': RANDOM_STATE
        }

    def _setup_directories(self):
        self.config['output_dir'].mkdir(parents=True, exist_ok=True)
        (self.config['output_dir'] / 'figures').mkdir(exist_ok=True)
        (self.config['output_dir'] / 'models').mkdir(exist_ok=True)
        (self.config['output_dir'] / 'reports').mkdir(exist_ok=True)

    def load_and_process_data(self):
        df = pd.read_excel(self.config['input_file'],
                           sheet_name=self.config['sheet_name'])
        print(f"数据加载完成: {len(df)} 个样本")
        abnormal_col = self.config['abnormal_col']
        if abnormal_col in df.columns:
            # 直接将染色体的非整倍体列作为标签（0=正常，1=异常）
            df['label'] = pd.to_numeric(df[abnormal_col], errors='coerce').fillna(0).astype(int)
            # 统计
            n_normal = (df['label'] == 0).sum()
            n_abnormal = (df['label'] == 1).sum()
            n_pregnant = df[self.config['pregnant_id_col']].nunique()

            df.to_excel(self.config['output_dir'] / 'processed_data.xlsx',
                        index=False)
        else:
            raise ValueError(f"未找到标签列: {abnormal_col}")

        self.data = df
        return df

    def prepare_features(self, df):
        exclude_cols = ['label', '异常标签', self.config['pregnant_id_col'],
                        self.config['abnormal_col'], self.config['weight_col']]
        numeric_cols = []
        for col in df.columns:
            if col not in exclude_cols:
                try:
                    numeric_data = pd.to_numeric(df[col], errors='coerce')
                    if numeric_data.notna().mean() > 0.3:
                        numeric_cols.append(col)
                except Exception:
                    pass
        print(f"\n自动识别 {len(numeric_cols)} 个数值特征")
        
        X = df[numeric_cols].copy()
        for col in X.columns:
            X[col] = pd.to_numeric(X[col], errors='coerce')
        X = X.dropna(axis=1, how='all')
        valid_cols = X.columns.tolist()
        print(f"有效特征: {len(valid_cols)} 个")

        imputer = SimpleImputer(strategy='median')
        X_imputed = imputer.fit_transform(X)
        X = pd.DataFrame(X_imputed, columns=valid_cols, index=df.index)

        self.feature_names = valid_cols
        y = df['label'].values

        if self.config['weight_col'] in df.columns:
            weights = pd.to_numeric(df[self.config['weight_col']], errors='coerce').fillna(1.0).values
            print(f"✓ 使用样本权重列: {self.config['weight_col']}")
            print(f"  权重范围: [{weights.min():.3f}, {weights.max():.3f}]")
            print(f"  平均权重: {weights.mean():.3f}")
        else:
            print(" 未找到权重列，使用默认权重1.0")
            weights = np.ones(len(y))

        groups = df[self.config['pregnant_id_col']].values
        print(f"✓ 使用孕妇分组列: {self.config['pregnant_id_col']}")

        unique, counts = np.unique(y, return_counts=True)
        print(f"✓ 标签分布: {dict(zip(unique, counts))}")
        return X, y, weights, groups

    def split_by_pregnant(self, X, y, weights, groups, test_size=0.3):
        print("\n按孕妇分组切分数据...")
        unique_groups = np.unique(groups)
        group_labels = []
        for group_id in unique_groups:
            group_mask = groups == group_id
            group_y = y[group_mask]
            group_label = int(np.median(group_y))
            group_labels.append(group_label)
        group_labels = np.array(group_labels)

        train_groups, test_groups = train_test_split(
            unique_groups, test_size=test_size,
            random_state=self.config['random_state'],
            stratify=group_labels
        )

        train_mask = np.isin(groups, train_groups)
        test_mask = np.isin(groups, test_groups)

        X_train = X[train_mask]; X_test = X[test_mask]
        y_train = y[train_mask]; y_test = y[test_mask]
        weights_train = weights[train_mask]; weights_test = weights[test_mask]
        groups_train = groups[train_mask]; groups_test = groups[test_mask]

        print(f"  训练集: {len(X_train)} 样本, {len(train_groups)} 位孕妇")
        print(f"  测试集: {len(X_test)} 样本, {len(test_groups)} 位孕妇")
        print(f"  训练集异常率: {(y_train==1).mean()*100:.2f}%")
        print(f"  测试集异常率: {(y_test==1).mean()*100:.2f}%")

        overlap = set(groups_train) & set(groups_test)
        if len(overlap) > 0:
            print(f"⚠️ 警告：有 {len(overlap)} 位孕妇同时出现在训练集和测试集！")
        else:
            print("✓ 确认：没有孕妇同时出现在训练集和测试集")

        return (X_train, y_train, weights_train, groups_train,
                X_test, y_test, weights_test, groups_test)

    def create_resampled_datasets(self, X, y, weights, groups):
        print("\n" + "="*60)
        print("步骤2: 生成重采样数据集")
        print("="*60)

        datasets = {'original': (X, y, weights, groups)}

        unique, counts = np.unique(y, return_counts=True)
        if len(unique) < 2:
            print("⚠️ 数据只有一个类别，无法进行重采样")
            return datasets

        min_class_count = min(counts)
        print(f"✓ 少数类样本数: {min_class_count}")

        resamplers = {}
        if min_class_count >= 2:
            resamplers['random_oversample'] = RandomOverSampler(random_state=RANDOM_STATE)
            resamplers['random_undersample'] = RandomUnderSampler(random_state=RANDOM_STATE)
        if min_class_count >= 6:
            k_neighbors = min(5, min_class_count - 1)
            resamplers['smote'] = SMOTE(random_state=RANDOM_STATE, 
                                        k_neighbors=k_neighbors)
            resamplers['borderline_smote'] = BorderlineSMOTE(random_state=RANDOM_STATE, 
                                                             k_neighbors=k_neighbors)
            resamplers['smoteenn'] = SMOTEENN(random_state=RANDOM_STATE, 
                                              smote=SMOTE(k_neighbors=k_neighbors))

        for name, resampler in resamplers.items():
            try:
                if 'smote' in name.lower() or 'enn' in name.lower():
                    X_with_weight = X.copy()
                    X_with_weight['weight'] = weights
                    X_res, y_res = resampler.fit_resample(X_with_weight, y)
                    weights_res = X_res['weight'].values
                    X_res = X_res.drop(['weight'], axis=1)
                    n_resampled = len(y_res)
                    if 'enn' in name.lower():
                        groups_res = np.array([f'synthetic_{i}' for i in range(n_resampled)])
                    else:
                        n_original = len(y)
                        groups_res = np.array(['synthetic'] * n_resampled)
                        if n_resampled >= n_original:
                            groups_res[:n_original] = groups
                else:
                    indices = np.arange(len(y))
                    indices_res, y_res = resampler.fit_resample(indices.reshape(-1, 1), y)
                    indices_res = indices_res.ravel()
                    X_res = X.iloc[indices_res].reset_index(drop=True)
                    weights_res = weights[indices_res]
                    groups_res = groups[indices_res]

                datasets[name] = (X_res, y_res, weights_res, groups_res)
                n_normal = int(np.sum(y_res == 0))
                n_abnormal = int(np.sum(y_res == 1))
                ratio = n_abnormal / n_normal if n_normal > 0 else 0
                print(f"✓ {name:20s}: 正常={n_normal:4d}, 异常={n_abnormal:4d}, 比例={ratio:.3f}")
            except Exception as e:
                print(f"✗ {name:20s}: 失败 - {str(e)}")

        stats_df = self._create_resample_stats(datasets)
        stats_df.to_excel(self.config['output_dir'] / 'resample_stats.xlsx', index=False)
        return datasets

    def _create_resample_stats(self, datasets):
        stats = []
        for name, data_tuple in datasets.items():
            y = data_tuple[1]
            n_normal = int(np.sum(y == 0))
            n_abnormal = int(np.sum(y == 1))
            total = len(y)
            pos_ratio = n_abnormal / total * 100 if total > 0 else 0
            class_ratio = n_abnormal / n_normal if n_normal > 0 else 0
            stats.append({
                'Dataset': name,
                'Normal': n_normal,
                'Abnormal': n_abnormal,
                'Total': total,
                'Positive_Rate': pos_ratio,
                'Class_Ratio': class_ratio
            })
        return pd.DataFrame(stats)

    def train_models(self, datasets):
        print("\n" + "="*60)
        print("步骤3: 模型训练与评估")
        print("="*60)

        model_configs = {
            'GradientBoosting': {
                'model': GradientBoostingClassifier(random_state=RANDOM_STATE),
                'param_grid': {
                    'n_estimators': [50, 100, 200],
                    'learning_rate': [0.01, 0.05, 0.1],
                    'max_depth': [3, 5, 7],
                    'min_samples_split': [10, 20],
                    'min_samples_leaf': [5, 10],
                    'subsample': [0.8, 1.0]
                }
            },
            'AdaBoost': {
                'model': AdaBoostClassifier(random_state=RANDOM_STATE, algorithm='SAMME'),
                'param_grid': {
                    'n_estimators': [50, 100, 200],
                    'learning_rate': [0.01, 0.05, 0.1, 0.5]
                }
            }
        }
        results = {}

        for data_name, data_tuple in datasets.items():
            print(f"\n处理数据集: {data_name}")
            print("-" * 40)

            X = data_tuple[0]
            y = data_tuple[1]
            weights = data_tuple[2] if len(data_tuple) > 2 else np.ones(len(y))
            groups = data_tuple[3] if len(data_tuple) > 3 else np.arange(len(y))

            (X_train, y_train, weights_train, groups_train,
             X_test, y_test, weights_test, groups_test) = self.split_by_pregnant(
                X, y, weights, groups, test_size=self.config['test_ratio']
            )

            scaler = StandardScaler()
            X_train_scaled = scaler.fit_transform(X_train)
            X_test_scaled = scaler.transform(X_test)

            data_results = {}

            for model_name, config in model_configs.items():
                print(f"  训练 {model_name}...")

                n_unique_groups = len(np.unique(groups_train))
                n_splits = min(self.config['cv_folds'], n_unique_groups)

                cv = GroupKFold(n_splits=n_splits)
                groups_for_cv = groups_train

                if len(y_train) < 100:
                    simple_grid = {}
                    for key, values in config['param_grid'].items():
                        if key == 'n_estimators':
                            simple_grid[key] = [50, 100]
                        elif key == 'learning_rate':
                            simple_grid[key] = [0.05, 0.1]
                        elif key == 'max_depth':
                            simple_grid[key] = [3, 5]
                        else:
                            simple_grid[key] = values[:2] if len(values) > 1 else values
                    param_grid = simple_grid
                else:
                    param_grid = config['param_grid']

                def weighted_auc_scorer(estimator, X_, y_):
                    y_pred_ = estimator.predict_proba(X_)[:, 1]
                    return roc_auc_score(y_, y_pred_)

                grid = GridSearchCV(
                    config['model'], param_grid,
                    cv=cv if groups_for_cv is None else cv.split(X_train_scaled, y_train,
                                                                 groups_for_cv),
                    scoring=weighted_auc_scorer,
                    n_jobs=-1, verbose=0
                )

                try:
                    grid.fit(X_train_scaled, y_train, sample_weight=weights_train)
                    best_model = grid.best_estimator_

                    train_metrics = self._evaluate_model(best_model, X_train_scaled, 
                                                         y_train, weights_train)
                    test_metrics = self._evaluate_model(best_model, X_test_scaled, 
                                                        y_test, weights_test)

                    y_test_proba = best_model.predict_proba(X_test_scaled)[:, 1]

                    data_results[model_name] = {
                        'model': best_model,
                        'scaler': scaler,
                        'train_metrics': train_metrics,
                        'test_metrics': test_metrics,
                        'best_params': grid.best_params_,
                        'cv_score': grid.best_score_,
                        'X_test': X_test_scaled,
                        'y_test': y_test,
                        'y_test_proba': y_test_proba,
                        'weights_test': weights_test,
                        'groups_test': groups_test
                    }

                    print(f"    CV AUC: {grid.best_score_:.4f}")
                    print(f"    测试 AUC: {test_metrics['auc']:.4f}, "
                          f"F1: {test_metrics['f1']:.4f}, "
                          f"召回率: {test_metrics['recall']:.4f}")

                except Exception as e:
                    print(f"    训练失败: {str(e)}")

            results[data_name] = data_results

        self.results = results
        return results

    def _evaluate_model(self, model, X, y, sample_weight=None):
        y_pred = model.predict(X)
        y_proba = model.predict_proba(X)[:, 1]

        try:
            mcc = matthews_corrcoef(y, y_pred, sample_weight=sample_weight)
        except TypeError:
            mcc = matthews_corrcoef(y, y_pred)
        try:
            kappa = cohen_kappa_score(y, y_pred, sample_weight=sample_weight)
        except TypeError:
            kappa = cohen_kappa_score(y, y_pred)

        metrics = {
            'accuracy': accuracy_score(y, y_pred, sample_weight=sample_weight),
            'precision': precision_score(y, y_pred, sample_weight=sample_weight, zero_division=0),
            'recall': recall_score(y, y_pred, sample_weight=sample_weight),
            'f1': f1_score(y, y_pred, sample_weight=sample_weight),
            'auc': roc_auc_score(y, y_proba, sample_weight=sample_weight),
            'ap': average_precision_score(y, y_proba, sample_weight=sample_weight),
            'mcc': mcc,
            'kappa': kappa,
            'brier': brier_score_loss(y, y_proba, sample_weight=sample_weight)
        }

        cm = confusion_matrix(y, y_pred, sample_weight=sample_weight)
        if cm.shape == (2, 2):
            tn, fp, fn, tp = cm.ravel()
            metrics['specificity'] = tn / (tn + fp) if (tn + fp) > 0 else 0
            metrics['npv'] = tn / (tn + fn) if (tn + fn) > 0 else 0

        return metrics

    def plot_comparison(self):
        print("\n" + "="*60)
        print("步骤4: 生成可视化报告")
        print("="*60)

        self._plot_metrics_comparison()
        self._plot_roc_curves()
        self._plot_best_model_analysis()
        self._plot_weight_analysis()

        print("✓ 可视化报告已生成")

    def _plot_weight_analysis(self):
        if not hasattr(self, 'results') or not self.results:
            return

        fig, axes = plt.subplots(1, 2, figsize=(12, 5))

        best_result = None
        for data_name, models in self.results.items():
            for model_name, result in models.items():
                if 'weights_test' in result:
                    best_result = result
                    break
            if best_result:
                break

        if best_result and 'weights_test' in best_result:
            weights = best_result['weights_test']
            y_test = best_result['y_test']
            y_proba = best_result['y_test_proba']

            axes[0].hist(weights[y_test == 0], bins=20, alpha=0.5, label='正常', density=True)
            axes[0].hist(weights[y_test == 1], bins=20, alpha=0.5, label='异常', density=True)
            axes[0].set_xlabel('样本权重 (w_Q)')
            axes[0].set_ylabel('密度')
            axes[0].set_title('按类别的权重分布')
            axes[0].legend()
            axes[0].grid(True, alpha=0.3)

            scatter = axes[1].scatter(weights, y_proba, c=y_test, cmap='coolwarm', alpha=0.6)
            axes[1].set_xlabel('样本权重 (w_Q)')
            axes[1].set_ylabel('预测概率')
            axes[1].set_title('权重与预测的关系')
            plt.colorbar(scatter, ax=axes[1], label='真实标签')
            axes[1].grid(True, alpha=0.3)

        plt.suptitle('样本权重分析', fontsize=14, y=1.02)
        plt.tight_layout()
        plt.savefig(self.config['output_dir'] / 'figures' / 'weight_analysis.png',
                    dpi=200, bbox_inches='tight')
        plt.close()

    def _plot_metrics_comparison(self):
        import numpy as np
        import pandas as pd
        import matplotlib as mpl
        import matplotlib.pyplot as plt
        import matplotlib.ticker as mticker
        from matplotlib.patches import Patch

        mpl.rcParams['font.sans-serif'] = ['Microsoft YaHei', 'SimHei', 'Arial Unicode MS',
                                           'DejaVu Sans']
        mpl.rcParams['axes.unicode_minus'] = False
        mpl.rcParams.update({
        'axes.facecolor': 'white',
        'figure.facecolor': 'white',
        'axes.edgecolor': '#E0E0E0',
        'axes.linewidth': 1.0,
        'axes.titlesize': 12,
        'axes.titleweight': 'bold',
        'axes.labelsize': 11,
        'xtick.labelsize': 10,
        'ytick.labelsize': 10,
        'legend.fontsize': 10,
        'grid.color': '#E6E6E6',
        'grid.linestyle': ':',
        'grid.linewidth': 0.8,
        })

        comparison_data = []
        for data_name, models in self.results.items():
            for model_name, result in models.items():
                comparison_data.append({
                'Dataset': data_name,
                'Model': model_name,
                'AUC': result['test_metrics']['auc'],
                'F1': result['test_metrics']['f1'],
                'Recall': result['test_metrics']['recall'],
                'Precision': result['test_metrics']['precision'],
                'MCC': result['test_metrics']['mcc']
            })
        df_comp = pd.DataFrame(comparison_data)

        # 固定数据集/模型顺序
        dataset_order = list(pd.unique(df_comp['Dataset']))
        model_order_raw = list(pd.unique(df_comp['Model']))
        df_comp['Dataset'] = pd.Categorical(df_comp['Dataset'], categories=dataset_order, ordered=True)

        display_model_order = [_zh_model(m) for m in model_order_raw]
        cmap = plt.get_cmap('tab10')
        color_map = {m_disp: cmap(i) for i, m_disp in enumerate(display_model_order)}
        hatches_seq = ['/', '\\\\', 'x', '.', '-', '+', 'o', 'O', '*']
        hatch_map = {m_disp: hatches_seq[i % len(hatches_seq)] for i, m_disp in enumerate(display_model_order)}

        fig, axes = plt.subplots(2, 3, figsize=(16, 10))
        axes = axes.flatten()
        metrics = ['AUC', 'F1', 'Recall', 'Precision', 'MCC']

        def annotate_bars(ax, fmt='{:.2f}'):
            for container in ax.containers:
                ax.bar_label(container, fmt=fmt, padding=2, fontsize=9)


        for idx, metric in enumerate(metrics):
            ax = axes[idx]
            pivot_df = df_comp.pivot(index='Dataset', columns='Model', values=metric)

            pivot_df = pivot_df.reindex(columns=model_order_raw)
            pivot_df.index = [_zh_dataset(i) for i in pivot_df.index]
            rename_map = {orig: disp for orig, disp in zip(model_order_raw, display_model_order)}
            pivot_df = pivot_df.rename(columns=rename_map)
            col_colors = [color_map[c] for c in pivot_df.columns]
            bars = pivot_df.plot(kind='bar', ax=ax, width=0.76,
                             color=col_colors, edgecolor='#444', linewidth=0.8, zorder=2)
            for j, container in enumerate(ax.containers):
                for p in container.patches:
                    p.set_hatch(hatch_map[pivot_df.columns[j]])
            ax.set_title(f'{_zh_metric(metric)}对比')
            ax.set_xlabel('')
            ax.grid(axis='y', alpha=0.6, zorder=0)
            for spine in ['top', 'right']:
                ax.spines[spine].set_visible(False)
            ax.set_xticklabels(ax.get_xticklabels(), rotation=18, ha='right')

            if metric == 'MCC':
                ax.set_ylim(-1.0, 1.0)
                ax.yaxis.set_major_locator(mticker.MultipleLocator(0.5))
            else:
                ax.set_ylim(0.0, 1.0)
                ax.yaxis.set_major_locator(mticker.MultipleLocator(0.2))
            ax.yaxis.set_major_formatter(mticker.FormatStrFormatter('%.1f'))
            annotate_bars(ax, fmt='{:.2f}')
            ax.legend_.remove()
        axes[-1].axis('off')
        legend_handles = [Patch(facecolor=color_map[m], edgecolor='#444',
                            hatch=hatch_map[m], label=m) for m in display_model_order]
        fig.legend(handles=legend_handles, loc='upper center',
               ncol=min(4, len(legend_handles)), frameon=True, fancybox=True,
               framealpha=0.9, borderpad=0.4, bbox_to_anchor=(0.5, 1.01))
        plt.suptitle('不同数据集的模型性能对比', fontsize=16, weight='bold', y=1.06)
        plt.tight_layout(rect=[0.02, 0.02, 0.98, 0.98])
        out_dir = self.config['output_dir'] / 'figures'
        out_dir.mkdir(parents=True, exist_ok=True)
        plt.savefig(out_dir / 'metrics_comparison.png', dpi=300, bbox_inches='tight')
        plt.close()

    def _plot_roc_curves(self):
        fig, axes = plt.subplots(2, 3, figsize=(15, 10))
        dataset_names = list(self.results.keys())
        for idx, data_name in enumerate(dataset_names[:6]):
            ax = axes[idx // 3, idx % 3]
            if data_name in self.results:
                for model_name, result in self.results[data_name].items():
                    y_test = result['y_test']
                    y_proba = result['y_test_proba']
                    weights_test = result.get('weights_test', None)
                    fpr, tpr, _ = roc_curve(y_test, y_proba, sample_weight=weights_test)
                    auc = roc_auc_score(y_test, y_proba, sample_weight=weights_test)
                    ax.plot(fpr, tpr, label=f'{_zh_model(model_name)} (AUC={auc:.3f})', lw=2)
                ax.plot([0, 1], [0, 1], 'k--', lw=1, alpha=0.5)
                ax.set_xlabel('假阳性率')
                ax.set_ylabel('真阳性率')
                ax.set_title(f'ROC曲线 - {_zh_dataset(data_name)}')
                ax.legend(loc='lower right')
                ax.grid(True, alpha=0.3)
        for idx in range(len(dataset_names), 6):
            fig.delaxes(axes[idx // 3, idx % 3])

        plt.suptitle('不同数据集（加权）的ROC曲线', fontsize=14, y=1.02)
        plt.tight_layout()
        plt.savefig(self.config['output_dir'] / 'figures' / 'roc_curves.png',
                    dpi=200, bbox_inches='tight')
        plt.close()

\end{lstlisting}
\end{appendices}

